\section{System: platform and mission protocol}\label{sec:system}

\subsection{Six-repository research ecosystem}
\sysname{} spans six co-developed repositories: TensorLBM
(lattice-Boltzmann fluid dynamics in PyTorch), TensorFVM (finite-volume
flows), TensorFEM (structural analysis), TensorDEM (discrete-element
media), TensorPhysics (scenario orchestration, validation vocabulary,
run manifests, artifact stores), and TensorPlatform (the FastAPI/Vue
platform hosting the mission system). A shared workspace lock governs
local checkouts; every scenario run records the code revisions of all
participating repositories in its run manifest.

\subsection{Mission protocol}
A \emph{mission} is a durable task object with a goal, a primary
repository and base ref, optional \emph{co-repositories} for
cross-project composition, a budget triple (tool calls, wall seconds,
tokens), a pinned model, and an origin (owner or agent). Its life cycle:

\begin{itemize}
\item \textbf{Claim.} Resident daemons poll a claim endpoint; a single
SQL \texttt{UPDATE\ ...\ WHERE id = (SELECT ... FOR UPDATE SKIP
LOCKED)} atomically moves the oldest eligible queued mission to
\texttt{claimed} and mints a one-time claim token. Claims whose
heartbeat went silent for more than 90\,s are lazily reaped as
\texttt{agent\_lost}.
\item \textbf{Execute.} The claiming daemon prepares jailed per-mission
git worktrees (one per repository in $\{repo\} \cup co\_repos$, branch
\texttt{agent/m<id>}, push disabled via per-worktree \texttt{pushurl}),
then drives an OpenHands SDK conversation with a sanitized environment
whitelist. Bridge tools (\S\ref{sec:method}) expose commit, test runs,
literature search, URL fetch with domain and private-IP guards,
acceptance requests, and orchestration.
\item \textbf{Heartbeat.} A watcher thread reports progress (tool count,
elapsed seconds, HEAD commit) every 15\,s and enforces the budget; a
watchdog escalates a wedged pause by process exit, relying on the flock
wrapper for restart and the reaper for cleanup.
\item \textbf{Terminal acceptance.} An agent may only finish by
requesting acceptance, which parks the mission in
\texttt{awaiting\_acceptance} and opens an approval row; an administrator
(owner) approves or rejects. Approval decisions never rest with agents.
All agent--platform interactions append immutable events (plan, tool
call/result, commit, web search, approval, decision) to a per-session
event stream.
\end{itemize}

\subsection{Model-pinned daemon fleet}
Seven relay-accessible backbone models each run a dedicated daemon
behind a per-model lock. The claim endpoint filters the queue by the
daemon's pinned model; exactly one designated daemon additionally claims
legacy missions with no model pin. Execution always uses the mission's
own pinned model, so claim routing and execution model are separable
concerns. Roles (\S\ref{sec:method}) resolve to models through an
owner-editable mapping file re-read on every dispatch.

\subsection{Research-run console}
The same platform operates a scenario execution service with the same
claim/reap discipline for solver runs (the stage ladder of
\S\ref{sec:case}), so \emph{development} missions and \emph{research}
executions share one audit substrate: append-only events, approval
chains, artifact checksums, and executor heartbeats. This unification is
what later allows the fleet's behavior to be studied with the same rigor
as the physics it produces (\S\ref{sec:ablation}).
